\subsection{Morfismos de presentación finita.}
\label{subsection:IV.1.6-fr}

\begin{definition}[1.6.1]
\label{IV.1.6.1-fr}
Sean $X$, $Y$ dos preesquemas. Se dice que un morfismo $f:X\to Y$ es
de presentación finita si verifica las tres condiciones siguientes:
\begin{enumerate}
  \item[\textup{(i)}] $f$ es localmente de presentación finita;
  \item[\textup{(ii)}] $f$ es cuasicompacto \textup{(lo que, cuando se verifica (i),
  equivale a decir que $f$ es \emph{de tipo finito} (1.5.2))};
  \item[\textup{(iii)}] $f$ es cuasiseparado.
\end{enumerate}
Se dice entonces también que $X$ es de presentación finita sobre $Y$, o un
$Y$-preesquema de presentación finita.
\end{definition}

La condición (iii) se verifica automáticamente si $f$ es separado, o si $X$
es localmente noetheriano (1.2.8); cuando $Y$ es localmente noetheriano,
es lo mismo decir que $f$ es \emph{de presentación finita} o que $f$
es \emph{de tipo finito} (esta última condición implica que $X$ es
localmente noetheriano (1.3.3)).

\oldpage[IV-1]{235}

\begin{proposition}[1.6.2]
\label{IV.1.6.2-fr}
\begin{enumerate}
  \item[\textup{(i)}] Toda inmersión cuasicompacta y localmente de presentación
  finita (en particular toda inmersión abierta cuasicompacta) es de
  presentación finita.
  \item[\textup{(ii)}] El compuesto de dos morfismos de presentación finita es de
  presentación finita.
  \item[\textup{(iii)}] Sean $X$, $Y$ dos $S$-preesquemas, $f:X\to Y$ un
  $S$-morfismo de presentación finita. Entonces, para todo morfismo $S'\to S$,
  el morfismo $f_{(S')}:X_{(S')}\to Y_{(S')}$ es de presentación finita.
  \item[\textup{(iv)}] Sean $f:X\to Y$, $f':X'\to Y'$ dos $S$-morfismos de
  presentación finita; entonces
  $f\times_Sf':X\times_SX'\to Y\times_SY'$ es de presentación finita.
  \item[\textup{(v)}] Si el compuesto $g\circ f$ de dos morfismos es de presentación
  finita y si $g$ es cuasiseparado y localmente de presentación finita, $f$
  es de presentación finita.
\end{enumerate}
\end{proposition}

Esto resulta inmediatamente de (1.1.2), (1.2.2), (1.2.4) y (1.4.3).

Resulta en particular de (1.6.2, (iii)) que si $f$ es un morfismo de
presentación finita y $U$ un abierto de $Y$, la restricción
$f^{-1}(U)\to U$ de $f$ es un morfismo de presentación finita. Recíprocamente,
sea $(U_\alpha)$ un recubrimiento de $Y$ por abiertos afines, y supongamos
que las restricciones $f^{-1}(U_\alpha)\to U_\alpha$ de $f$ son
morfismos de presentación finita; resulta entonces que $f$ es de
presentación finita, pues $f$ es evidentemente localmente de presentación finita,
cuasicompacto en virtud de (1.1.1) y cuasiseparado en virtud de (1.2.6). Se puede,
por tanto, decir que la propiedad de un morfismo $f:X\to Y$ de ser de presentación
finita es \emph{local sobre $Y$}.

Si $X$ es un preesquema cuasiseparado, todo morfismo $f:X\to Y$ es
cuasiseparado (1.2.3). Por tanto, si $f$ es cuasicompacto y localmente de
presentación finita, $f$ es entonces de presentación finita.

En particular:

\begin{corollary}[1.6.3]
\label{IV.1.6.3-fr}
Sea $\varphi:A\to B$ un homomorfismo de anillos. Para que el morfismo
correspondiente $f:\operatorname{Spec}(B)\to\operatorname{Spec}(A)$ sea de presentación finita, es necesario y
suficiente que $\varphi$ haga de $B$ una $A$-álgebra de presentación finita.
\end{corollary}

La necesidad de la condición resulta, en efecto, de (1.4.6).

\begin{remark}[1.6.4]
\label{IV.1.6.4-fr}
En la definición (1.6.1), la condición (iii) no es consecuencia de las otras
dos. Sea, por ejemplo, $Y$ un esquema afín cuyo espacio subyacente no
sea noetheriano, y sea $U$ un abierto no cuasicompacto en $Y$. Sea
$X$ el preesquema obtenido pegando dos preesquemas $Y_1$, $Y_2$ isomorfos a
$Y$ a lo largo de los abiertos $U_1$, $U_2$ correspondientes a $U$ ($\mathbf{0}_{\mathrm I}$,
4.1.7), de modo que $X$ es unión de dos abiertos isomorfos respectivamente
a $Y_1$ y $Y_2$ (y que identificamos con estos últimos), con
$Y_1\cap Y_2=U$. Además, hay un morfismo $f:X\to Y$ que coincide en
$Y_i$ con el isomorfismo dado $Y_i\to Y$ ($i=1,2$). Está claro que $f$
verifica la condición (i) de (1.6.1), por ser un isomorfismo local; verifica
también (ii), pues la imagen inversa de un abierto cuasicompacto de $Y$ es la
unión de sus imágenes en $Y_1$ y $Y_2$; pero como $Y_1\cap Y_2$ no es
cuasicompacto, $f$ no es cuasiseparado (1.2.7).
\end{remark}

\begin{proposition}[1.6.5]
\label{IV.1.6.5-fr}
Sean $X'$, $X''$ dos preesquemas, $X=X'\amalg X''$ su suma,
$f':X'\to Y$, $f'':X''\to Y$ dos morfismos, $f:X\to Y$ el morfismo que
coincide con $f'$ en $X'$ y con $f''$ en $X''$. Para que $f$ sea de
presentación finita, es necesario y suficiente que $f'$ y $f''$ lo sean.
\end{proposition}

\par\medskip\noindent
\begin{proof}
Basta mostrar que para que $f$ posea una de las tres propiedades de la
definición
\oldpage[IV-1]{236}
(1.6.1), es necesario y suficiente que $f'$ y $f''$ la posean; es evidente
para la propiedad de ser localmente de presentación finita, que es local sobre
$X$; para la propiedad de ser cuasicompacto, se vio en (1.1.6) y para
la propiedad de ser cuasiseparado, en (1.2.9).
\end{proof}

